\chapter{Operators, Spectra and Banach Algebras}

\begin{center}\itshape
A matrix is understood through its eigenvalues. The chapters above were that theory for matrices: invariant subspaces, inner products, the spectral theorem, singular values, graphs, inequalities, and linear dynamics. An operator on an infinite-dimensional space may have none at all, and yet it has a
spectrum, never empty, whose size is read off the norms of its powers. Behind this lies an algebra, and behind the algebra, a space.
\end{center}

In 1932 Norbert Wiener proved a lemma that looks innocent. If a function on the circle has an absolutely convergent Fourier series and
never vanishes, then $\frac1f$ also has an absolutely convergent Fourier series. For a trigonometric polynomial with no zeros, the claim
is not even obvious: $\frac1{2+\cos t}$ is smooth, so its coefficients decay fast, but what about a merely absolutely convergent
series, whose sum may be nowhere differentiable? Wiener's proof was a few pages of delicate estimates, and he needed the lemma for his
great Tauberian theorem (Chapter~19).

Nine years later, a twenty-eight-year-old in Moscow gave a proof in three lines. Israel Gelfand had arrived in the capital at sixteen
without finishing school, worked as a doorman at the Lenin Library and as a tutor, and attended university lectures without being
enrolled, until Kolmogorov took him as a graduate student at nineteen. His 1941 theory of ``normed rings'' looked at the functions with
absolutely convergent Fourier series as an \emph{algebra}, asked what its maximal ideals are, and found the points of the circle. Once
the algebra is understood, Wiener's lemma says only this: an element that no point of the circle ``sees'' as zero is invertible. The
proof became the advertisement for a whole way of thinking, in which analysis is done by algebra and geometry is recovered from the
algebra of functions on it, a philosophy that runs through Grothendieck's algebraic geometry and Connes's noncommutative geometry.

This chapter first studies the operators that behave most like matrices, the compact ones, for which Fredholm's alternative recovers
the dichotomy of linear algebra: an equation $x-Kx=y$ either has a unique solution for every $y$, or the homogeneous equation has
solutions. Then it defines the spectrum of an operator, proves that it is never empty, and computes its size by Gelfand's
spectral radius formula. Gelfand's theory of commutative Banach algebras follows, and Wiener's lemma falls out.

Throughout, $H$ is a Hilbert space and $E$ a Banach space over $\C$ unless said otherwise; complex scalars are needed for spectra, as
the rotation of the plane by a right angle, which has no real eigenvalue, already shows.

\section{Adjoints}

For $T\in\mathcal L(H)=\mathcal L(H,H)$ and fixed $y$, $x\mapsto\langle Tx,y\rangle$ is a bounded linear form, so by the Riesz--Fr\'echet theorem (Theorem~2.3) there is a unique $T^*y$ with $\langle Tx,y\rangle=\langle x,T^*y\rangle$ for all
$x$. The map $T^*$, the \emph{adjoint} of $T$, is linear. For matrices it is the conjugate transpose.

\begin{proposition}\label{prop:adjoint}
$T^{**}=T$, $(ST)^*=T^*S^*$, $\|T^*\|=\|T\|$ and $\|T^*T\|=\|T\|^2$. Moreover $\ker T^*=(\operatorname{ran}T)^\perp$ and $\overline{\operatorname{ran}T}=(\ker T^*)^\perp$.
\end{proposition}
\begin{proof}
The first two identities follow from the definition and uniqueness. $\|T^*y\|^2=\langle TT^*y,y\rangle\le\|T\|\|T^*y\|\|y\|$, so $\|T^*\|\le\|T\|$, and applying this to $T^*$ gives equality. Then
$\|T^*T\|\le\|T\|^2$, and $\|Tx\|^2=\langle T^*Tx,x\rangle\le\|T^*T\|\|x\|^2$ gives the reverse. Finally $y\in\ker T^*$ iff $\langle x,T^*y\rangle=0$ for all $x$ iff $y\perp\operatorname{ran}T$; take orthogonals and use
Corollary~2.2.
\end{proof}

$T$ is \emph{self-adjoint} if $T^*=T$, \emph{normal} if $T^*T=TT^*$, and \emph{unitary} if $T^*T=TT^*=I$. Orthogonal projections are self-adjoint (Corollary~2.2), multiplication by a bounded real function on $L^2$ is self-adjoint, and
the shift $S(x_1,x_2,\dots)=(0,x_1,x_2,\dots)$ on $\ell^2$ has adjoint $S^*(x_1,x_2,\dots)=(x_2,x_3,\dots)$: $S^*S=I$ but $SS^*\ne I$. An isometry that is not onto is something a matrix cannot be.

\section{Compact operators}

An operator $K:E\to F$ between Banach spaces is \emph{compact} if the image of the unit ball is relatively compact, that is, if every bounded sequence $(x_n)$ has a subsequence with $(Kx_n)$ convergent. Operators of finite rank are compact
(bounded sets of a finite-dimensional space are relatively compact), the identity of an infinite-dimensional space is not (Theorem~1.3), and compact operators are the infinite-dimensional operators closest to matrices.

\begin{proposition}\label{prop:compactops}
(a) The compact operators form a closed subspace of $\mathcal L(E,F)$, and $AKB$ is compact when $K$ is compact and $A,B$ bounded. (b) On a Hilbert space, an operator is compact if and only if it is a norm limit of operators of finite rank, and
$K$ is compact if and only if $K^*$ is.
\end{proposition}
\begin{proof}
(a) Sums and products with bounded operators are clear from the sequential definition. Let $\|K_n-K\|\to0$ with $K_n$ compact. A set is relatively compact in the complete space $F$ iff it is totally bounded (Volume~6, Theorem~4.8). Given
$\varepsilon$, pick $n$ with $\|K_n-K\|<\varepsilon$ and a finite $\varepsilon$-net $y_1,\dots,y_m$ of $K_n(B)$; then the $y_i$ form a $2\varepsilon$-net of $K(B)$.

(b) Let $K$ be compact on $H$ and $\varepsilon>0$. Take an $\varepsilon$-net $y_1,\dots,y_m$ of $K(B)$ and $P$ the orthogonal projection on their span. For $\|x\|\le1$ pick $i$ with $\|Kx-y_i\|<\varepsilon$; as $Py_i=y_i$ and $\|P\|\le1$,
$\|Kx-PKx\|\le\|Kx-y_i\|+\|P(y_i-Kx)\|<2\varepsilon$. So $\|K-PK\|\le2\varepsilon$, and $PK$ has finite rank. Conversely, limits of finite-rank operators are compact by (a). If $K=\lim F_n$ with $F_n$ of finite rank, then
$K^*=\lim F_n^*$ (Proposition~\ref{prop:adjoint}: $\|K^*-F_n^*\|=\|K-F_n\|$), and $F_n^*$ has finite rank, since $\ker F_n^*=(\operatorname{ran}F_n)^\perp$ has finite codimension. Apply this to $K^*$ for the converse.
\end{proof}

The Banach-space version of (b) is false: Enflo's goose (Chapter~1) is precisely a compact operator that is not a limit of finite-rank operators. The statement about adjoints survives (Schauder's theorem, Exercise~\ref{ex:schauder}).

\begin{example}[integral operators]\label{ex:hilbertschmidt}
Let $(X,\mu)$ be $\sigma$-finite and $k\in L^2(\mu\otimes\mu)$. Then $(Tf)(x)=\int k(x,y)f(y)\,d\mu(y)$ defines a compact operator on $L^2(\mu)$ with $\|T\|\le\|k\|_{L^2}$. Indeed, Cauchy--Schwarz and Tonelli give
$\int|Tf(x)|^2\,d\mu(x)\le\int\!\!\int|k(x,y)|^2\,d\mu(y)\,d\mu(x)\cdot\|f\|_2^2$. Functions $\sum_{i\le m}c_i\mathbf 1_{A_i}(x)\mathbf 1_{B_i}(y)$ with $A_i,B_i$ of finite measure are dense in $L^2(\mu\otimes\mu)$ (they include the indicators of
finite unions of rectangles of finite measure, which approximate every product-measurable set of finite measure, by the construction of the product measure, Volume~6, Theorem~8.4); the corresponding operators have rank at most
$m$, and $\|T-T_m\|\le\|k-k_m\|_{L^2}$. By Proposition~\ref{prop:compactops}, $T$ is compact. These are the \emph{Hilbert--Schmidt} integral operators, the kernels of Fredholm's equations.
\end{example}

\section{The Fredholm alternative}

\begin{theorem}[Fredholm alternative]\label{thm:fredholm}
Let $K$ be a compact operator on a Hilbert space $H$. Then:
\begin{enumerate}
\item[(a)] $\ker(I-K)$ is finite-dimensional, and $\operatorname{ran}(I-K)$ is closed, equal to $\ker(I-K^*)^\perp$;
\item[(b)] $I-K$ is injective if and only if it is surjective;
\item[(c)] $\dim\ker(I-K)=\dim\ker(I-K^*)$.
\end{enumerate}
So either $x-Kx=y$ has a unique solution for every $y$, or the homogeneous equation $x=Kx$ has non-zero solutions; and then $x-Kx=y$ is solvable exactly when $y$ is orthogonal to the finitely many independent solutions of $z=K^*z$.
\end{theorem}
\begin{proof}
Write $N=\ker(I-K)$. (a) On $N$, $K$ is the identity, so the unit ball of $N$ is $K(B\cap N)$, relatively compact and closed: $N$ is finite-dimensional by Theorem~1.3. Let $y_n=x_n-Kx_n\to y$; replacing $x_n$ by $x_n-P_Nx_n$, we may
assume $x_n\perp N$. If $(x_n)$ were unbounded, then along a subsequence $\|x_n\|\to\infty$, and $z_n=x_n/\|x_n\|$ would satisfy $z_n-Kz_n\to0$; a further subsequence has $Kz_n\to w$, hence $z_n\to w$, with $\|w\|=1$,
$w-Kw=0$ and $w\perp N$: absurd. So $(x_n)$ is bounded, a subsequence has $Kx_n\to v$, then $x_n=y_n+Kx_n\to y+v=:x$ and $x-Kx=y$. The range is closed, and equals $\ker(I-K^*)^\perp$ by Proposition~\ref{prop:adjoint}.

(b) Suppose $I-K$ injective but $H_1=\operatorname{ran}(I-K)\ne H$. Let $H_n=(I-K)^nH$. Since $(I-K)^n=I-K_n$ with $K_n$ compact, each $H_n$ is closed by (a). As $(I-K)^n$ is injective and maps $H$ onto $H_n$ and $H_1$ onto
$H_{n+1}$, $H_{n+1}\ne H_n$. Pick unit vectors $e_n\in H_n\cap H_{n+1}^\perp$. For $m<n$,
\[Ke_m-Ke_n=e_m-\bigl(e_n+(I-K)e_m-(I-K)e_n\bigr),\]
and the bracket lies in $H_{m+1}$, which is orthogonal to $e_m$; so $\|Ke_m-Ke_n\|\ge1$, and $(Ke_n)$ has no convergent subsequence, contradicting compactness. Conversely, if $I-K$ is surjective, then
$\ker(I-K^*)=\operatorname{ran}(I-K)^\perp=\{0\}$; by the first half applied to $K^*$, $I-K^*$ is surjective, so $\ker(I-K)=\operatorname{ran}(I-K^*)^\perp=\{0\}$.

(c) Let $d=\dim N$ and $d^*=\dim\ker(I-K^*)$, and suppose $d<d^*$. Let $A:N\to\ker(I-K^*)$ be linear, injective and not onto, and $K'=K+AP_N$, compact. If $x-K'x=0$, then $x-Kx=AP_Nx$ lies in $\operatorname{ran}(I-K)$ and in
$\ker(I-K^*)=\operatorname{ran}(I-K)^\perp$, so both vanish: $P_Nx=0$ and $x\in N$, hence $x=0$. By (b), $I-K'$ is onto. But a non-zero $y\in\ker(I-K^*)$ orthogonal to $A(N)$ is orthogonal to $\operatorname{ran}(I-K)$ and to $\operatorname{ran}A$,
hence to $\operatorname{ran}(I-K')$: a contradiction. So $d\ge d^*$, and by symmetry ($K^{**}=K$) $d=d^*$.
\end{proof}

\begin{beauty}
For a square matrix $M$, ``$Mx=y$ is solvable for every $y$'' and ``$Mx=0$ only for $x=0$'' are the same statement, because the rank theorem counts dimensions. In infinite dimension the count is meaningless, the shift
is injective and not onto, and yet for $I-K$ the dichotomy survives. The compact perturbation $K$ is ``small'' enough that $I-K$ still behaves like a matrix. Fredholm found this in 1903 for integral equations, before
Hilbert spaces existed; it is the reason why the Dirichlet problem, and many other boundary value problems, have solutions.
\end{beauty}

\section{The spectrum}

A \emph{Banach algebra} is a Banach space $A$ with an associative bilinear product satisfying $\|ab\|\le\|a\|\|b\|$ and a unit $1$ with $\|1\|=1$. The examples to keep in mind are $\mathcal L(E)$ with composition, $C(K)$ with the
pointwise product, and $\ell^1(\Z)$ with convolution $(a*b)_n=\sum_ka_kb_{n-k}$, for which $\|a*b\|_1\le\|a\|_1\|b\|_1$ by Tonelli for series, with unit $\delta_0$.

\begin{lemma}[Neumann series]\label{lem:neumann}
If $\|a\|<1$, then $1-a$ is invertible with $(1-a)^{-1}=\sum_{n\ge0}a^n$ and $\|(1-a)^{-1}\|\le\frac1{1-\|a\|}$. The invertible elements form an open set $G$, on which $a\mapsto a^{-1}$ is continuous.
\end{lemma}
\begin{proof}
The series converges absolutely since $\|a^n\|\le\|a\|^n$, and $(1-a)\sum_{n\le N}a^n=1-a^{N+1}\to1$, and the same on the other side. If $b$ is invertible and $\|c\|<\|b^{-1}\|^{-1}$, then $b-c=b(1-b^{-1}c)$ is invertible with
$(b-c)^{-1}=\sum_n(b^{-1}c)^nb^{-1}$, and $\|(b-c)^{-1}-b^{-1}\|\le\sum_{n\ge1}\|b^{-1}\|^{n+1}\|c\|^n\to0$ as $c\to0$.
\end{proof}

The \emph{spectrum} of $a\in A$ is $\sigma(a)=\{\lambda\in\C:\lambda-a\text{ is not invertible}\}$ (we write $\lambda$ for $\lambda1$), and its \emph{spectral radius} is $r(a)=\sup\{|\lambda|:\lambda\in\sigma(a)\}$. For a matrix, the
spectrum is the set of eigenvalues. For an operator $T$, every eigenvalue is in $\sigma(T)$, but $\lambda-T$ can fail to be invertible while being injective: it may fail to be onto.

\begin{theorem}[Gelfand]\label{thm:spectrum}
For every $a$ in a Banach algebra, $\sigma(a)$ is a non-empty compact subset of $\{|\lambda|\le\|a\|\}$, and
\[r(a)=\lim_{n\to\infty}\|a^n\|^{1/n}=\inf_n\|a^n\|^{1/n}.\]
\end{theorem}
\begin{proof}
If $|\lambda|>\|a\|$, then $\lambda-a=\lambda(1-a/\lambda)$ is invertible by Lemma~\ref{lem:neumann}, with
\[(\lambda-a)^{-1}=\sum_{n\ge0}\frac{a^n}{\lambda^{n+1}}.\tag{$*$}\]
The complement of $\sigma(a)$ is open by the same lemma, so $\sigma(a)$ is compact. For $\varphi\in A^*$, the function $f_\varphi(\lambda)=\varphi((\lambda-a)^{-1})$ is holomorphic on $\C\setminus\sigma(a)$: near $\lambda_0$,
$(\lambda-a)^{-1}=\sum_n(\lambda_0-\lambda)^n(\lambda_0-a)^{-n-1}$ (the proof of Lemma~\ref{lem:neumann} with $b=\lambda_0-a$ and $c=(\lambda_0-\lambda)$), a power series in $\lambda-\lambda_0$.

\emph{Non-empty.} If $\sigma(a)=\varnothing$, every $f_\varphi$ is entire, and by $(*)$, $|f_\varphi(\lambda)|\le\frac{\|\varphi\|}{|\lambda|-\|a\|}\to0$ as $|\lambda|\to\infty$. By Liouville's theorem (Volume~4, Chapter~9) $f_\varphi=0$, so
$\varphi(a^{-1})=f_\varphi(0)=0$ for all $\varphi$, and $a^{-1}=0$ by Hahn--Banach (Corollary~1.5): absurd, since $aa^{-1}=1\ne0$.

\emph{The formula.} If $\lambda\in\sigma(a)$, then $\lambda^n\in\sigma(a^n)$, because $\lambda^n-a^n=(\lambda-a)(\lambda^{n-1}+\cdots+a^{n-1})$ with commuting factors, so if $\lambda^n-a^n$ were invertible, so would be $\lambda-a$. Hence
$|\lambda|^n\le\|a^n\|$, and $r(a)\le\inf_n\|a^n\|^{1/n}$. Conversely, $f_\varphi$ is holomorphic on $|\lambda|>r(a)$, and equal there to $\sum_n\varphi(a^n)\lambda^{-n-1}$ for $|\lambda|>\|a\|$. By uniqueness of Laurent expansions on the
annulus $\{|\lambda|>r(a)\}$ (Volume~4, Chapter~9), the series converges for all $|\lambda|>r(a)$. So for such $\lambda$, $\varphi(a^n/\lambda^n)\to0$ for every $\varphi$, and by the uniform boundedness principle
(Theorem~1.7, applied to the forms $J(a^n/\lambda^n)$ on $A^*$) $\|a^n\|\le C_\lambda|\lambda|^n$. Thus $\limsup\|a^n\|^{1/n}\le|\lambda|$ for every $|\lambda|>r(a)$, and $\limsup\|a^n\|^{1/n}\le r(a)\le\inf\|a^n\|^{1/n}$.
\end{proof}

The spectral radius depends only on the algebra structure and the norm of powers, not on the particular norm: equivalent norms give the same $r(a)$. For a matrix it is the largest modulus of an eigenvalue, and the formula says
that $\|M^n\|$ grows like $r(M)^n$ up to subexponential factors, the rate at which iterative methods converge (Exercise~\ref{ex:jacobi}).

\begin{example}\label{ex:volterraex}
The \emph{Volterra operator} $(Vf)(x)=\int_0^xf$ on $L^2([0,1])$ (or $C([0,1])$) satisfies $|V^nf(x)|\le\frac{x^{n-1}}{(n-1)!}\|f\|_1$, by induction, so $\|V^n\|\le\frac1{(n-1)!}$ and $r(V)=0$: $\sigma(V)=\{0\}$. But $V$ has no eigenvalue at
all: $Vf=\lambda f$ with $\lambda\ne0$ gives $f$ continuous, then $C^1$, with $f=\lambda f'$ and $f(0)=0$, so $f=0$; and $Vf=0$ gives $\int_0^xf=0$ for all $x$, so $f=0$ a.e.\ (Volume~6, Theorem~10.9). The spectrum $\{0\}$ is
there only because $V$ is not onto.
\end{example}

\begin{figure}[ht]
\centering
\fig{ch03-pseudospectra}
\caption{The $64\times64$ Jordan block $J$ (ones above the diagonal) has only the eigenvalue $0$. But $\|(z-J)^{-1}\|$ exceeds $10^{2}$, $10^{6}$, $10^{10}$, $10^{14}$ inside the circles shown, which crowd towards the unit circle: as the
size grows, matrices whose spectrum is $\{0\}$ ``look like'' the shift, whose spectrum is the closed unit disc (Exercise~\ref{ex:shift}).}
\label{fig:pseudospectra}
\end{figure}

\begin{think}
Can an operator have no eigenvalue at all? Example~\ref{ex:volterraex} says yes, and so does the shift $S$ on $\ell^2$. Can its spectrum be empty? Theorem~\ref{thm:spectrum} says never, and the proof needs complex analysis. Over $\R$,
the rotation by a right angle on $\R^2$ has empty real spectrum: the theorem is truly about $\C$.
\end{think}

\section{Gelfand's theory of commutative Banach algebras}

From now on $A$ is a \emph{commutative} Banach algebra with unit, over $\C$. A \emph{character} of $A$ is a non-zero algebra homomorphism $\chi:A\to\C$; then $\chi(1)=1$. Gelfand's idea is to view the characters as the points of a space, and
each $a\in A$ as a function on it, $\hat a(\chi)=\chi(a)$.

\begin{lemma}\label{lem:gelfandmazur}
(a) (Gelfand--Mazur) A Banach algebra in which every non-zero element is invertible is $\C1$. (b) Every character satisfies $|\chi(a)|\le\|a\|$, hence is continuous with norm $1$. (c) If $I$ is a closed proper ideal, $A/I$ with the
norm $\|a+I\|=\inf_{i\in I}\|a+i\|$ is a Banach algebra. (d) Maximal ideals are closed.
\end{lemma}
\begin{proof}
(a) By Theorem~\ref{thm:spectrum} there is $\lambda\in\sigma(a)$; $\lambda-a$ is not invertible, so it is $0$. (b) If $|\chi(a)|>\|a\|$, then $\chi(a)-a$ is invertible (Lemma~\ref{lem:neumann}), but a homomorphism maps invertible elements to
non-zero numbers, while $\chi(\chi(a)-a)=0$. (c) The quotient norm is a norm because $I$ is closed, and it is submultiplicative ($(a+i)(b+j)\in ab+I$). If $\|i+1\|<1$ for some $i\in I$, then $-i=1-(1+i)$ would be invertible, and $I=A$;
so $\|1+I\|=1$. Completeness: given a Cauchy sequence, pass to a subsequence with $\|(a_{n+1}-a_n)+I\|<2^{-n}$, choose representatives $d_n$ of $a_{n+1}-a_n$ with $\|d_n\|<2^{-n}$; then $a_1+\sum_nd_n$ converges in $A$, and its class is
the limit. (d) A proper ideal contains no invertible element, and the invertibles form an open set containing the ball $\|1-x\|<1$; so the closure of a proper ideal stays at distance $1$ from $1$ and is proper. A maximal ideal equals its closure.
\end{proof}

\begin{theorem}[Gelfand, 1941]\label{thm:gelfand}
Let $\Delta$ be the set of characters of $A$, with the weak* topology of $A^*$.
\begin{enumerate}
\item[(a)] $\chi\mapsto\ker\chi$ is a bijection from $\Delta$ onto the set of maximal ideals of $A$.
\item[(b)] $a$ is invertible if and only if $\chi(a)\ne0$ for every $\chi\in\Delta$.
\item[(c)] $\sigma(a)=\{\chi(a):\chi\in\Delta\}$, and $r(a)=\max_{\chi\in\Delta}|\chi(a)|$.
\item[(d)] $\Delta$ is a compact Hausdorff space, and the \emph{Gelfand transform} $a\mapsto\hat a$ is an algebra homomorphism $A\to C(\Delta)$ with $\|\hat a\|_\infty=r(a)\le\|a\|$.
\end{enumerate}
\end{theorem}
\begin{proof}
(a) $\ker\chi$ has codimension $1$ ($a-\chi(a)1\in\ker\chi$), so it is maximal. Conversely, a maximal ideal $M$ is closed (Lemma~\ref{lem:gelfandmazur}(d)), so $A/M$ is a Banach algebra; it is a field, because for $a\notin M$ the ideal
$M+aA$ is $A$, so $ab\in1+M$ for some $b$. By Gelfand--Mazur, $A/M=\C1$, and $a\mapsto$ (the scalar of $a+M$) is a character with kernel $M$. A character is determined by its kernel, since $\chi(a)$ is the unique $\lambda$ with $\lambda-a\in\ker\chi$.

(b) If $a$ is invertible, $\chi(a)\chi(a^{-1})=1$. If not, $aA$ is a proper ideal ($1\notin aA$ as $A$ is commutative). The union of a chain of proper ideals is a proper ideal (none contains $1$), so by Zorn's lemma $aA$ is contained in
a maximal ideal $\ker\chi$, and $\chi(a)=0$. (c) Apply (b) to $\lambda-a$. The maximum is attained because $\sigma(a)$ is compact.

(d) $\Delta$ lies in the unit ball of $A^*$ (Lemma~\ref{lem:gelfandmazur}(b)), which is weak* compact (Theorem~1.16), and $\Delta$ is weak* closed: it is the set of $\varphi$ in the ball with $\varphi(ab)=\varphi(a)\varphi(b)$ for all $a,b$ and
$\varphi(1)=1$, each condition being closed. The weak* topology is Hausdorff, and each $\hat a$ is continuous by definition of that topology. $\widehat{ab}=\hat a\hat b$ since characters are multiplicative, and $\|\hat a\|_\infty=r(a)$ by (c).
\end{proof}

For $A=C(K)$ with $K$ compact metric, the characters are the evaluations $f\mapsto f(x)$ (Exercise~\ref{ex:ckcharacters}), $\Delta=K$, and the Gelfand transform is the identity: the space is recovered from the algebra of its functions.

\section{Wiener's lemma}

Let $W$ be the algebra of functions $f(t)=\sum_{n\in\Z}a_ne^{int}$ on $\mathbb T$ with $\|f\|_W=\sum_n|a_n|<\infty$. The map $a\mapsto f$ identifies $\ell^1(\Z)$ with convolution and $W$ with the pointwise product: the product of two
absolutely convergent series is the series of the convolution, by Fubini for series (Volume~0). So $W$ is a commutative Banach algebra with unit $1=e^{i0t}$.

\begin{proposition}\label{prop:wcharacters}
The characters of $W$ are the evaluations $f\mapsto f(t_0)$, $t_0\in\mathbb T$.
\end{proposition}
\begin{proof}
Evaluations are characters. Conversely, let $\chi$ be a character and $z=\chi(e^{it})$. As $\|e^{it}\|_W=\|e^{-it}\|_W=1$, Lemma~\ref{lem:gelfandmazur}(b) gives $|z|\le1$ and $|z^{-1}|=|\chi(e^{-it})|\le1$, so $z=e^{it_0}$. Then
$\chi(e^{int})=z^n=e^{int_0}$ for all $n\in\Z$, and since $f=\sum_na_ne^{int}$ converges in $W$ and $\chi$ is continuous, $\chi(f)=\sum_na_ne^{int_0}=f(t_0)$.
\end{proof}

\begin{theorem}[Wiener's lemma, 1932]\label{thm:wiener}
If $f\in W$ does not vanish on $\mathbb T$, then $\frac1f\in W$: the reciprocal of a nowhere-zero absolutely convergent Fourier series is an absolutely convergent Fourier series.
\end{theorem}
\begin{proof}
By Proposition~\ref{prop:wcharacters}, $\chi(f)\ne0$ for every character $\chi$, so $f$ is invertible in $W$ by Theorem~\ref{thm:gelfand}(b): there is $g\in W$ with $fg=1$.
\end{proof}

\begin{figure}[ht]
\centering
\fig{ch03-wiener}
\caption{The curve $t\mapsto f(t)=1+0.6e^{it}+0.3e^{-3it}$, $t\in[0,2\pi]$, avoids the origin. Wiener's lemma then guarantees that $\frac1f$ has absolutely summable Fourier coefficients. The proof never computes them: it only checks
that no character of $W$, that is, no point of the circle, sends $f$ to $0$.}
\label{fig:wiener}
\end{figure}

The Gelfand transform of $W$ is the identity map, and that of $\ell^1(\Z)$ is the Fourier series. The analogue for $L^1(\R)$ with convolution, whose characters are given by the Fourier transform, is Wiener's Tauberian theorem
(Chapter~19). The same three lines prove much more: if $F$ is holomorphic near the range of $f\in W$, then $F\circ f\in W$ (Wiener--L\'evy), by the holomorphic functional calculus.

\begin{think}
Why does Wiener's lemma need the whole space of characters? Could an element $f$ of some algebra be non-zero at every point of a set $X$ of characters and still fail to be invertible? (Try the disc algebra of functions continuous on
the closed disc and holomorphic inside, with $X$ the unit circle and $f(z)=z$.)
\end{think}

\section{Exercises}

\begin{exercise}[$\star$]\label{ex:volterraeq}
Show that for every $\lambda\in\C$ and $g\in C([0,1])$ the equation $f(x)-\lambda\int_0^xf(t)\,dt=g(x)$ has a unique solution $f\in C([0,1])$, given by a series. Solve it for $g=1$.
\end{exercise}

\begin{exercise}[$\star$]\label{ex:fredholmrank1}
Solve $u(x)-\lambda\int_0^1xy\,u(y)\,dy=g(x)$ in $L^2([0,1])$ for each $\lambda\in\C$, and check the Fredholm alternative on this example.
\end{exercise}

\begin{exercise}[$\star\star$]\label{ex:shift}
Let $S$ be the shift on $\ell^2$. (a) Show that $S$ has no eigenvalue and that every $\lambda$ with $|\lambda|<1$ is an eigenvalue of $S^*$. (b) Show that $\sigma(S)=\sigma(S^*)$ is the closed unit disc. (c) Let $J$ be the $N\times N$ matrix
with ones above the diagonal. Show that $\sigma(J)=\{0\}$ but $\|(z-J)^{-1}\|\ge|z|^{-N}$ for $0<|z|<1$.
\end{exercise}

\begin{exercise}[$\star\star$]\label{ex:multiplication}
Let $\mu$ be $\sigma$-finite, $\varphi\in L^\infty(\mu)$ and $M_\varphi f=\varphi f$ on $L^2(\mu)$. Show that $\|M_\varphi\|=\|\varphi\|_\infty$, that $\sigma(M_\varphi)$ is the \emph{essential range} $\{\lambda:\mu(|\varphi-\lambda|<\varepsilon)>0$ for all
$\varepsilon>0\}$, and that $\lambda$ is an eigenvalue iff $\mu(\varphi=\lambda)>0$. What is the spectrum of multiplication by $x$ on $L^2([0,1])$?
\end{exercise}

\begin{exercise}[$\star$]\label{ex:jacobi}
\emph{Iterative solvers.} Let $M=D-N$ be a complex $n\times n$ matrix with $D$ its diagonal part, and $B=D^{-1}N$. Show that the Jacobi iteration $x_{k+1}=D^{-1}(Nx_k+b)$ converges to the solution of $Mx=b$ for every $x_0$ if and only if
$r(B)<1$, and that this holds when $M$ is strictly diagonally dominant ($|m_{ii}|>\sum_{j\ne i}|m_{ij}|$ for all $i$).
\end{exercise}

\begin{exercise}[$\star\star$]\label{ex:jacobson}
In a Banach algebra, show that if $1-xy$ is invertible then so is $1-yx$, with inverse $1+y(1-xy)^{-1}x$. Deduce that $\sigma(ab)\cup\{0\}=\sigma(ba)\cup\{0\}$ and $r(ab)=r(ba)$. Find $a,b$ with $\sigma(ab)\ne\sigma(ba)$.
\end{exercise}

\begin{exercise}[$\star$]\label{ex:ckcharacters}
Let $K$ be a compact metric space. Show that every character of $C(K)$ is an evaluation $f\mapsto f(x)$ for a unique $x\in K$.
\end{exercise}

\begin{exercise}[$\star\star\star$]\label{ex:schauder}
\emph{Schauder's theorem.} Let $K:E\to F$ be a compact operator between Banach spaces. Show that $K^*:F^*\to E^*$, $K^*\varphi=\varphi\circ K$, is compact.
\end{exercise}

\section{Hints}

\begin{hint}{ex:volterraeq}
$\|V^n\|\le\frac1{(n-1)!}$, so $\sum_n\lambda^nV^n$ converges in norm whatever $\lambda$.
\end{hint}
\begin{hint}{ex:fredholmrank1}
Put $c=\int_0^1yu(y)\,dy$; then $u=g+\lambda cx$. Find the equation satisfied by $c$.
\end{hint}
\begin{hint}{ex:shift}
(a) Try $x=(1,\lambda,\lambda^2,\dots)$. (b) The spectrum is closed, and $\lambda-T$ is invertible iff $\overline\lambda-T^*$ is. (c) Expand $(z-J)^{-1}=\sum_{k<N}J^kz^{-k-1}$ and look at one entry.
\end{hint}
\begin{hint}{ex:multiplication}
If $\lambda$ is outside the essential range, $\frac1{\varphi-\lambda}$ is bounded. If inside, test $M_\varphi-\lambda$ on indicators of subsets of $\{|\varphi-\lambda|<\varepsilon\}$.
\end{hint}
\begin{hint}{ex:jacobi}
The error $e_k=x_k-x$ satisfies $e_{k+1}=Be_k$. Use Theorem~\ref{thm:spectrum}, and the max-row-sum norm for diagonal dominance.
\end{hint}
\begin{hint}{ex:jacobson}
Multiply out $(1-yx)(1+ycx)$ with $c=(1-xy)^{-1}$. For the example, use the shift.
\end{hint}
\begin{hint}{ex:ckcharacters}
If no $x$ works, for each $x$ find $f_x\in\ker\chi$ with $f_x(x)\ne0$; use compactness to build $\sum|f_{x_i}|^2\in\ker\chi$, nowhere zero.
\end{hint}
\begin{hint}{ex:schauder}
The closure $X$ of $K(B_E)$ is compact; a bounded sequence in $F^*$, restricted to $X$, is equicontinuous. Use Arzel\`a--Ascoli.
\end{hint}

\section{Solutions}

\begin{solution}{ex:volterraeq}
\idea The Volterra operator is ``smaller than every multiple of the identity'' in the sense of spectral radius, so $I-\lambda V$ is always invertible by a Neumann series, whatever the size of $\lambda$.
\writeup By Example~\ref{ex:volterraex} (the same induction works on $C([0,1])$ with the sup norm), $\|\lambda^nV^n\|\le\frac{|\lambda|^n}{(n-1)!}$, so $R=\sum_{n\ge0}\lambda^nV^n$ converges in norm, and $(I-\lambda V)R=R(I-\lambda V)=I$ by
telescoping. So $f=Rg$ is the unique solution. For $g=1$, $V^n1=\frac{x^n}{n!}$, so $f=\sum_n\frac{(\lambda x)^n}{n!}=e^{\lambda x}$; check: $e^{\lambda x}-\lambda\int_0^xe^{\lambda t}\,dt=1$.
\end{solution}

\begin{solution}{ex:fredholmrank1}
\idea A kernel of rank one turns the integral equation into one scalar equation for one unknown number, and the alternative is visible: the scalar equation is solvable uniquely unless its coefficient vanishes.
\writeup With $c=\int_0^1yu(y)\,dy$, the equation says $u=g+\lambda cx$. Multiplying by $x$ and integrating, $c=\int_0^1yg+\frac{\lambda c}3$, that is $(1-\frac\lambda3)c=\int_0^1yg$. If $\lambda\ne3$, then $c=\frac{3\int yg}{3-\lambda}$ and
$u=g+\frac{3\lambda\int yg}{3-\lambda}x$ is the unique solution. If $\lambda=3$, a solution exists iff $\int_0^1yg(y)\,dy=0$, and then $u=g+3cx$ for any $c$. Here $K=3\langle\cdot,x\rangle x$ is compact and self-adjoint, $\ker(I-K)=\ker(I-K^*)=\C x$, and the
condition $\int yg=0$ is $g\perp\ker(I-K^*)$, as Theorem~\ref{thm:fredholm} predicts.
\end{solution}

\begin{solution}{ex:shift}
\idea The shift is the simplest operator that is not a matrix: an isometry onto a proper subspace. Its adjoint has a whole disc of eigenvectors, geometric sequences, and closedness of the spectrum fills in the boundary.
Truncating the shift kills all this spectrum except $0$, but the resolvent remembers it.
\writeup (a) If $Sx=\lambda x$, comparing coordinates, $\lambda x_1=0$ and $\lambda x_{n+1}=x_n$; if $\lambda=0$ then $x=0$ as $S$ is injective, and if $\lambda\ne0$ then $x_1=0$ and inductively all $x_n=0$. For $|\lambda|<1$, $x=(\lambda^{n-1})_n\in\ell^2$ and
$S^*x=(\lambda,\lambda^2,\dots)=\lambda x$. (b) $\|S\|=\|S^*\|=1$, so both spectra lie in the closed disc. $\sigma(S^*)$ contains the open disc by (a), and is closed, so it is the closed disc. Since $(\lambda-T)^*=\overline\lambda-T^*$
and an operator is invertible iff its adjoint is ($(T^{-1})^*=(T^*)^{-1}$), $\sigma(S)=\overline{\sigma(S^*)}$, the closed disc. (c) $J^N=0$, so $\lambda-J$ is invertible for $\lambda\ne0$ with $(z-J)^{-1}=\sum_{k<N}J^kz^{-k-1}$, while $J$ is not
invertible: $\sigma(J)=\{0\}$. The $(1,N)$ entry of $J^{N-1}$ is $1$ and that of $J^k$ is $0$ for $k<N-1$, so the $(1,N)$ entry of $(z-J)^{-1}$ is $z^{-N}$, and the norm of a matrix is at least the modulus of each entry.
\end{solution}

\begin{solution}{ex:multiplication}
\idea A multiplication operator is a diagonal matrix with a continuum of diagonal entries, the values of $\varphi$. Its spectrum is the set of values, up to null sets, and the eigenvalues are the values taken on a set of positive measure.
\writeup $\|\varphi f\|_2\le\|\varphi\|_\infty\|f\|_2$. For $\varepsilon>0$, the set $\{|\varphi|>\|\varphi\|_\infty-\varepsilon\}$ has positive measure and contains a set $B$ of finite positive measure; $\|\varphi\mathbf 1_B\|_2\ge(\|\varphi\|_\infty-\varepsilon)\|\mathbf 1_B\|_2$.
So $\|M_\varphi\|=\|\varphi\|_\infty$. If $\lambda$ is not in the essential range, $|\varphi-\lambda|\ge\varepsilon$ a.e.\ for some $\varepsilon$, and $M_{1/(\varphi-\lambda)}$ is a bounded inverse of $M_\varphi-\lambda$. If $\lambda$ is in it, take $B$ of finite positive
measure inside $\{|\varphi-\lambda|<\varepsilon\}$: $\|(M_\varphi-\lambda)\mathbf 1_B\|\le\varepsilon\|\mathbf 1_B\|$, so an inverse would have norm at least $\frac1\varepsilon$ for every $\varepsilon$: impossible. Finally $(\varphi-\lambda)f=0$ a.e.\ iff $f=0$
a.e.\ off $\{\varphi=\lambda\}$, which allows $f\ne0$ iff that set has positive measure. For $\varphi(x)=x$ on $[0,1]$, the spectrum is $[0,1]$ and there is no eigenvalue.
\end{solution}

\begin{solution}{ex:jacobi}
\idea Iterations converge when powers of the iteration matrix go to zero, and the spectral radius formula says that this happens exactly when the spectral radius is below one, whatever norm one prefers.
\writeup $x=D^{-1}(Nx+b)$ for the solution, so $e_k=B^ke_0$. If $r(B)<1$, pick $r(B)<q<1$; by Theorem~\ref{thm:spectrum}, $\|B^k\|\le q^k$ for large $k$, so $e_k\to0$ for every $e_0$. If $r(B)\ge1$, let $v$ be an eigenvector with
$|\mu|\ge1$; with $x_0=x+v$, $e_k=\mu^kv\not\to0$. Strict diagonal dominance says $\sum_j|b_{ij}|=\sum_{j\ne i}\frac{|m_{ij}|}{|m_{ii}|}<1$ for each $i$, so in the norm $\|x\|_\infty$ the operator norm of $B$, the maximal row sum, is less than $1$, and
$r(B)\le\|B\|<1$.
\end{solution}

\begin{solution}{ex:jacobson}
\idea Formally $(1-yx)^{-1}=\sum(yx)^n=1+y\bigl(\sum(xy)^n\bigr)x$, which suggests the formula; then it can be checked without any convergence.
\writeup With $c=(1-xy)^{-1}$: $(1-yx)(1+ycx)=1-yx+ycx-yxycx=1-yx+y(1-xy)cx=1-yx+yx=1$, and similarly $(1+ycx)(1-yx)=1-yx+ycx-ycxyx=1-yx+yc(1-xy)x=1$. For $\lambda\ne0$, $\lambda-ab=\lambda(1-\frac a\lambda b)$ is invertible iff
$1-b\frac a\lambda$ is, iff $\lambda-ba$ is. So the spectra agree away from $0$, and so do the spectral radii. With the shift, $S^*S=I$ has spectrum $\{1\}$, while $SS^*$, the projection onto $\{x_1=0\}$, has spectrum $\{0,1\}$.
\end{solution}

\begin{solution}{ex:ckcharacters}
\idea A character that is not an evaluation would have, for each point, a function in its kernel that does not vanish there. Compactness glues finitely many of them into a function with no zero, which is invertible, and an
invertible element cannot be in the kernel of a character.
\writeup Suppose $\ker\chi$ is not of the form $\{f:f(x)=0\}$ for any $x$. Then $\ker\chi\not\subseteq\{f:f(x)=0\}$ (both are maximal ideals, by Theorem~\ref{thm:gelfand}(a) for the second), so for each $x$ there is $f_x\in\ker\chi$ with
$f_x(x)\ne0$, hence $|f_x|>0$ on a neighbourhood $U_x$. Cover $K$ by $U_{x_1},\dots,U_{x_m}$. Then $g=\sum_if_{x_i}\overline{f_{x_i}}$ is in $\ker\chi$ (an ideal) and $g>0$ on $K$, so $g$ is invertible in $C(K)$: contradiction. So
$\chi(f)=0$ whenever $f(x)=0$, for some $x$; since $f-f(x)$ vanishes at $x$, $\chi(f)=f(x)$. The point is unique because continuous functions separate points.
\end{solution}

\begin{solution}{ex:schauder}
\idea The adjoint only sees the linear forms on the compact set $X=\overline{K(B_E)}$, and on a compact set a bounded family of $1$-Lipschitz functions is relatively compact by Arzel\`a--Ascoli.
\writeup Let $(\varphi_n)$ be a sequence in the unit ball of $F^*$. On the compact metric space $X$, the functions $\varphi_n|_X$ are bounded by $\sup_X\|y\|$ and $1$-Lipschitz, so equicontinuous; by Arzel\`a--Ascoli (Volume~6, Theorem~4.15) a
subsequence converges uniformly on $X$. Then $\|K^*\varphi_n-K^*\varphi_m\|=\sup_{\|x\|\le1}|(\varphi_n-\varphi_m)(Kx)|\le\sup_X|\varphi_n-\varphi_m|$, so $(K^*\varphi_n)$ is Cauchy along the subsequence, and converges in the
Banach space $E^*$.
\end{solution}

\section*{Sources and further reading}

\begin{itemize}
\item I.~Fredholm, ``Sur une classe d'\'equations fonctionnelles'', \emph{Acta Math.} 27 (1903), 365--390.
\item I.~M.~Gelfand, ``Normierte Ringe'', \emph{Mat. Sbornik} 9 (1941), 3--24; I.~M.~Gelfand, D.~A.~Raikov and G.~E.~Shilov, \emph{Commutative Normed Rings}, Chelsea, 1964.
\item N.~Wiener, ``Tauberian theorems'', \emph{Ann. of Math.} 33 (1932), 1--100. Lemma~IIe is Wiener's lemma.
\item F.~Riesz and B.~Sz.-Nagy, \emph{Functional Analysis}, Ungar, 1955. Compact operators and the Riesz--Schauder theory.
\item H.~Brezis, \emph{Functional Analysis}, Springer, 2011, Chapter~6; W.~Rudin, \emph{Functional Analysis}, 2nd ed., McGraw--Hill, 1991, Chapters~10--11.
\item L.~N.~Trefethen and M.~Embree, \emph{Spectra and Pseudospectra}, Princeton, 2005. Figure~\ref{fig:pseudospectra} and much more.
\item V.~Retakh (ed.), \emph{Israel Moiseevich Gelfand}, Notices Amer. Math. Soc. 60 (2013), 24--49 and 162--171. Memories of Gelfand's early years.
\end{itemize}
